%
% 4_RothsteinTrager.tex -- Abschnitt 4 Resultante anwendung in RT
%
% (c) 2026 Jakob Gierer, OST Ostschweizer Fachhochschule
%
% !TEX root = ../../buch.tex
% !TEX encoding = UTF-8
%
\section{Der Satz von Rothstein-Trager
\label{resultante:section:RT}}
\kopfrechts{Der Satz von Rothstein-Trager}
In der Anlaysis wird für die Berechnung von Stammfunktionen bei Integralen die Partialbruchzerlegung verwendet.
Diese Methode hat das Problem, dass die Nullstellen des Nennerpolynoms bekannt sein müssen.
Das Problem ist also die Residuen des Integranden zu berechnen ohne den Nenner zu faktorisieren.
Die Mathematiker Rothstein und Trager haben unabhängig voneinander den Satz \ref{buch:ratint:rothsteintrager:satz:rothsteintrager} bewiesen.
Der rationale Teil \(a(z)\) muss teilerfremd zu \(b(z)\) sein, mit \(\deg a(z)<\deg b(z)\).
Zusätzlich dazu muss das Polynom \(b(z)\) quadratfrei sein, dann sagt das Liouville-Prinzip
\begin{equation*}
    \int\frac{a(z)}{b(z)}\,dz=\sum_{k=1}^{m}c_k\log v_k(z),
    \label{resultante:equation:satzRT}
\end{equation*}
diese Form vorraus. Gesucht sind nun die \(c_k\) für
\begin{equation}
    a(z)-c_k\cdot b'(z) \quad \text{und} \quad b(z)
    \label{resultante:equation:polynome}
\end{equation}
die einen nitchttrivialen grössten gemeinsamen Teiler (Teiler \(\neq1\))
\begin{equation}
    v_k(z)=\ggt\left(a(z)-c_k\cdot b'(z),b(z)\right)
    \label{resultante:equation:vks}
\end{equation}
haben.

\begin{satz}
    \emph{(Rothstein-Trager)}. Die Nullstellen von \(\operatorname{R}(z)\) sind genau die Residuen von \(\frac{a(z)}{b(z)}\).
    \label{resultante:satz:RT} 
\end{satz}
\noindent
Die Nullstellen von \(\operatorname{R}(z)=\Res_z\) sind die gesuchten \(c_k\).

Dafür eignet sich die Resultante nach Definition \ref{resultante:definition:resultante} hervorragend.
Sie liefert als Lösung direkt den \(\ggt\) zweier Polynome mit dazugehöriger Nullstelle.

\subsection{Anwendung der Resultante
\label{resultante:subsection:anwendung}}
Für die Bestimmung von den Nullstellen \(c_k\) und den grössten gemeinsamen Teiler \(v_k(z)\) wird die Resultante
\begin{equation}
    \Res_z\left(a(z)-c_k\cdot b'(z),b(z)\right)=0
    \label{resultante:equation:resz}
\end{equation}
benötigt.
Dabei bedeutet das tiefgestellt \(z\), dass die Resultante nach \(z\) aufgestellt werden muss.
Dadurch wird die Variable \(z\), wie in Beispiel \ref{resultante:beispiel:elimination} aus der Resultanten \(\operatorname{R}(z)\) eliminiert.

\begin{beispiel}
    Die Stammfunktion von
    \begin{equation*}
        \int\frac{1}{z^3+z}\,dz
    \end{equation*}
    soll mit der Methoden von Rothstein-Trager und der Resultante bestimmt werden.
    Die Bedingungen für die Rothstein-Trager-Methode sind bereits erfüllt, mit \(a(z)=1\) und \(b(z)=z^3+z\).
    Es müssen daher die \(c_k\) und die dazugehörigen \(v_k(z)\) gefunden werden.
    Für die Resultante braucht es
    \begin{equation}
        \begin{aligned}
            a(z)&=1\\
            b(z)&=z^3+z\\
            b'(z)&=3z^2+1
        \end{aligned},
        \label{resultante:equation:factoren}
    \end{equation}
    sie sieht wie folgt aus:
    \begin{equation}
        \Res_z\left(-3c_kz^2+1-c_k,z^3+z\right)=0,
        \label{resultante:equation:resZgut}
    \end{equation}
    die daraus folgende Sylvester-Matrix ist (zur Vereinfachung \(c\) statt \(c_k\)):
    \begin{equation}
        \begin{NiceArray}{*{5}{c}}[cell-space-limits=3pt]
        -3c & 0 & 1-c & 0 & 0 \\
        0 & -3c & 0 & 1-c & 0 \\
        0 & 0 & -3c & 0 & 1-c \\
        1 & 0 & 1 & 0 & 0 \\
        0 & 1 & 0 & 1 & 0
        \CodeAfter
            \tikz {
                \draw ([xshift=-3pt,yshift=1pt]1-|1) rectangle ([xshift=4pt,yshift=-1pt]6-|6);
                \draw[dashed] ([xshift=-3pt,yshift=1pt]4-|1) -- ([xshift=4pt,yshift=1pt]4-|6);
            }
        \end{NiceArray}
        \,\,\to\,\,
        \begin{NiceArray}{*{5}{c}}[cell-space-limits=3pt]
        1 & 0 & \frac{c-1}{3c} & 0 & 0 \\
        0 & 1 & 0 & \frac{c-1}{3c} & 0 \\
        0 & 0 & 1 & 0 & \frac{c-1}{3c} \\
        1 & 0 & 1 & 0 & 0 \\
        0 & 1 & 0 & 1 & 0
        \CodeAfter
            \tikz {
                \draw ([xshift=-4pt,yshift=1pt]1-|1) rectangle ([xshift=3pt,yshift=-1pt]6-|6);
                \draw[dashed] ([xshift=-4pt,yshift=0.5pt]4-|1) -- ([xshift=3pt,yshift=0.5pt]4-|6);
            }
        \end{NiceArray}\,\,\,.
        \label{resultante:equation:sylGutStep1}
    \end{equation}
    Nach Anwendung der Zeilenoperationen erhält man das erste Tableau, auf dem kleineren Bereich kann nochmals eine Zeilenoperation durchgeführt werden:
    \begin{equation}
        \begin{NiceArray}{*{5}{c}}[cell-space-limits=3pt]
        1 & 0 & \frac{c-1}{3c} & 0 & 0 \\
        0 & 1 & 0 & \frac{c-1}{3c} & 0 \\
        0 & 0 & 1 & 0 & \frac{c-1}{3c} \\
        0 & 0 & \frac{2c+1}{3c} & 0 & 0 \\
        0 & 0 & 0 & \frac{2c+1}{3c} & 0
        \CodeAfter
            \tikz {
                \draw ([xshift=-4pt,yshift=1pt]1-|1) rectangle ([xshift=3pt,yshift=-1pt]6-|6);
                \draw[dashed] ([xshift=1pt]3-|3) -- ([xshift=1pt,yshift=-1pt]6-|3);
                \draw[dashed] ([xshift=1pt]3-|3) -- ([xshift=3pt]3-|6);
                \draw[dashed] ([xshift=1pt]4-|3) -- ([xshift=3pt]4-|6);
            }
        \end{NiceArray}
        \,\,\to\,\,
        \begin{NiceArray}{*{3}{c}}[cell-space-limits=3pt]
        1 & 0 & 0 \\
        0 & 1 & 0 \\
        0 & 0 & \frac{c-1}{3c}
        \CodeAfter
            \tikz {
                \draw ([xshift=-3pt,yshift=1pt]1-|1) rectangle ([xshift=3pt,yshift=-1pt]4-|4);
                \draw[dashed] ([xshift=-3pt]3-|1) -- ([xshift=3pt]3-|4);
            }
        \end{NiceArray}\,\,\,.
        \label{resultante:equation:sylGutStep2}
    \end{equation}
    Das Schlusstableau hat als entscheidendes Element:
    \begin{equation}
        \begin{aligned}
            \frac{c-1}{3c}&\overset{!}{=}0\\
            c&=1
        \end{aligned},
        \label{resultante:equation:solGut1}
    \end{equation}
    für \(c_1=1\) lässt sich im Schlusstableau aus (\ref{resultante:equation:sylGutStep2}) für \(v_1(z)\) den \(\ggt=z\) ablesen.
    Findige Leser konnten bereits im vorherigen Tableau aus (\ref{resultante:equation:sylGutStep2}) den Algorithmus beenden, nämlich mit:
    \begin{equation}
        \begin{aligned}
            \frac{2c+1}{3c}&\overset{!}{=}0\\
            2c+1&=0\\
            2c&=-1\\
            c&=-\frac{1}{2}
        \end{aligned},
        \label{resultante:equation:solGut2}
    \end{equation}
    für \(c_2=-\frac{1}{2}\) lässt sich ebenfalls im zweitletzten Tableau aus (\ref{resultante:equation:sylGutStep2}) für \(v_2(z)\) den \(\ggt=z^2+1\) ablesen.

    Setzt man die gefundenen Nullstellen zusammen mit ihren zugehörigen grössten gemeinsamen Teiler erhält man die Stammfunktion:
    \begin{equation*}
        \int\frac{1}{z^3+z}\,dz=c_1\log v_1(z)+c_2\log v_2(z)=1\log z-\frac{1}{2}\log(z^2+1).
    \end{equation*}
    \label{resultante:beispiel:gutesIntegral}
\end{beispiel}

\begin{beispiel}
    Die Methode soll mit einem weiteren Beispiel verdeutlicht werden.
    Das Integral
    \begin{equation*}
        \int\frac{1}{\log x}\,dx
    \end{equation*}
    soll gelöst werden.
    Dafür muss \(F(\vartheta)=\mathbb{Q}(x)\) mit dem logarithmischem Element \(\vartheta=\log x\) erweiterte werden.
    Der Integrand wird dadurch zu
    \begin{equation*}
        f=\frac{1}{\vartheta}
    \end{equation*}
    mit
    \begin{equation}
        \begin{aligned}
            a(\vartheta)&=1\\
            b(\vartheta)&=\vartheta\\
            Db(\vartheta)&=\frac{1}{x}
        \end{aligned}
    \end{equation}
    Das \(D\) steht hier für die Derivation nach \(x\).
    Die Resultante daraus sieht wie folgt aus:
    \begin{equation}
        \Res_\vartheta(a(\vartheta)-cDb(\vartheta),b(\vartheta))=\Res_\vartheta\left(1-c\frac{1}{x},\vartheta\right)=0,
        \label{resultante:equation:resTbad}
    \end{equation}
    der erste Polynom in der Resultante ist vom Grad \(0\) und das zweite vom Grad \(1\).
    Die daraus entstehende \(1\times1\)-Sylvester-Matrix im Tableau ist dann
    \begin{equation}
        \begin{NiceArray}{c}[cell-space-limits=3pt]
        1-c\frac{1}{x} \\
        \CodeAfter
            \tikz \draw ([xshift=-3pt,yshift=2pt]1-|1) rectangle ([xshift=3pt,yshift=-2pt]2-|2);
        \end{NiceArray}
        \,\,\,\,\,.
        \label{resultante:equation:sylBad}
    \end{equation}
    Der Algorithmus für dieses Tableau endet bevor er überhaupt das erste mal angewandt wird.
    Das kritische Element ist:
    \begin{equation}
        \begin{aligned}
            1-\frac{c}{x}&\overset{!}{=}0\\
            \frac{c}{x}&=1\\
            c&=x
        \end{aligned},
        \label{resultante:equation:solBad}
    \end{equation}
    die Nullstelle \(c\) muss für die Methode von Rothstein-Trager eine Konstante sein (\(Dc\overset{!}{=}0\)).
    Hier ist
    \begin{equation*}
        Dc=Dx=1\overset{!}{=}0 \,\text{\lightning}
    \end{equation*}
    Daraus folgt das Integral ist nicht elementar lösbar.
    Wie Beispiel \ref{buch:intalg:risch:bsp:1/logx} und Definition \ref{buch:intalg:risch:def:integrallogarithmus} zeigt, ist das der Integrallogarithmus.
    \label{resultante:beispiel:badIntegral}
\end{beispiel}

Sobald die Nullstellen \(c_k\) aus der Resultante keine Konstant sind, folgt daraus, dass solche Integrale nicht elementar lösbar sind.